\clearpage\label{answers}
{\sans\bfseries\LARGE\color{ink} Answers}\par
\footnotesize\linespread{0.94}\selectfont
Try every puzzle first. Where an answer shows a calculation, check the steps, not just the result.

\subhead{Chapter 1}
\textbf{Multix on the 8~cm.} No. 8~cm on the model stands for 400~cm in the room, because $8\times50=400$. A model is a stand-in, not the room. \textbf{Keep box.} Kept: size and place. Set aside: colour and how things feel.

\textbf{1.1} $200\div50=4$~cm.\\
\textbf{1.2} $3\times50=150$ and $2\times50=100$. So 150~cm by 100~cm.\\
\textbf{1.3} $1:20$. Each model centimetre stands for fewer real centimetres, so the model is bigger.\\
\textbf{1.4} (a) K (b) D (c) K (d) D (e) K.\\
\textbf{1.5} The bed and desk are at different scales, so their sizes do not match their real sizes.\\
\textbf{1.6} Chair 1~cm. Table 3~cm.\\
\textbf{1.7} $1.6\times50=80$~cm.\\
\textbf{1.8} $300\div50=6$.\\
\textbf{1.9} ``The model is a smaller stand-in that keeps the room's sizes and places.''\\
\textbf{1.10} Answers vary. Check that the scale is stated and that a kept and a dropped thing are named.

\subhead{Chapter 2}
\textbf{Multix on the half-size map.} No. Smaller fits the page but it drops more. A 50~cm window is 0.25~cm at 1~cm to 2~m, thinner than a pencil line. \textbf{Keep box.} Kept: door places and long lengths. Set aside: pictures and small things.

\textbf{2.1} $3\times2=6$~m.\\
\textbf{2.2} $8\div4=2$~cm.\\
\textbf{2.3} 0.5~cm and 0.25~cm.\\
\textbf{2.4} 1~cm to 10~cm. A room is a few metres, which fits that scale on a poster.\\
\textbf{2.5} (a) K (b) D (c) K (d) D.\\
\textbf{2.6} $2\times1.5=3$~cm. Yes, the bar was copied with the map, so it still says 3~cm is 4~m.\\
\textbf{2.7} The maps use different scales, so the lengths cannot be compared without the scales. Both hallways are 12~m.\\
\textbf{2.8} 6.\\
\textbf{2.9} ``The map keeps door places and lengths, and drops small things such as pictures.''\\
\textbf{2.10} Answers vary. Check the scale and at least one dropped thing.

\subhead{Chapter 3}
\textbf{Multix on every question.} No. A model answers only the questions its kept parts bear on. The door swing was not kept. \textbf{Keep box.} Kept: the question and the prediction. Set aside: anything the model was not built for.

\textbf{3.1} (a) M (b) S (c) D.\\
\textbf{3.2} $3.8\times50=190$~cm. No, the bed is longer than the gap.\\
\textbf{3.3} $150-120=30$~cm.\\
\textbf{3.4} No. Colour was not kept in a cardboard model.\\
\textbf{3.5} A model of one room cannot be tested against a different room.\\
\textbf{3.6} Chair $1.2\times50=60$~cm. Gap $1.0\times50=50$~cm. No, the chair is 10~cm too wide.\\
\textbf{3.7} Answers vary, for example whether the door opens fully.\\
\textbf{3.8} The answer circle says ``no''.\\
\textbf{3.9} ``The model says the bed does not fit, and it kept sizes but not door swing.''\\
\textbf{3.10} Cut a bike-sized outline at a scale and compare it with a shed outline at the same scale.

\subhead{Chapter 4}
\textbf{Multix on equal gaps.} No. The diagram keeps order and drops distance, so equal spacing means nothing about real distance. \textbf{Keep box.} Kept: order of stops, or distances. Set aside: distances, or the number of stops.

\textbf{4.1} $200+100+500=800$~m. $800\div100=8$~cm.\\
\textbf{4.2} (a) M (b) R (c) X.\\
\textbf{4.3} $500\div100=5$ times farther.\\
\textbf{4.4} Keeps order and count of stops. Drops distances.\\
\textbf{4.5} (a) M (b) R (c) M.\\
\textbf{4.6} The map, because walking depends on distance.\\
\textbf{4.7} The diagram, because counting stops depends on order.\\
\textbf{4.8} Map keeps distances. Diagram keeps order.\\
\textbf{4.9} ``Each diagram is better for a different question.''\\
\textbf{4.10} Answers vary. Check that each diagram has a kept and a dropped thing.

\subhead{Chapter 5}
\textbf{Multix on the 59~cm chair.} No. A 59~cm chair in a 60~cm gap has a margin of 1~cm, smaller than the model's biggest gap of 4~cm. \textbf{Keep box.} Kept: the prediction and the gap. Set aside: why it was off.

\textbf{5.1} $164-160=4$.\\
\textbf{5.2} $2.6\times50=130$~cm.\\
\textbf{5.3} $60-40=20$. Yes, 20 is bigger than 4.\\
\textbf{5.4} $60-59=1$. No, 1 is smaller than 4.\\
\textbf{5.5} Yes. The bed is 10~cm too long, and 10 is bigger than 4.\\
\textbf{5.6} One prediction could be lucky. Three show how big the gaps usually are.\\
\textbf{5.7} Different tools and cherry-picking one zero gap hide the real spread. Use one method and report every gap.\\
\textbf{5.8} Circles: 60, 58 and gap 2.\\
\textbf{5.9} ``The model's predictions were within 4~cm in the three tests we made.''\\
\textbf{5.10} Answers vary. Check that a prediction is written before measuring.

\subhead{Chapter 6}
\textbf{Multix on the bin.} No. The model is fine for large questions. Repair it for the small one and keep both, with the old model wrapped. \textbf{Keep box.} Kept: the old model and the fix. Set aside: the claim it replaced.

\textbf{Spot the cheat.} B. At 1~cm to 2~m its door is 1~cm, which would be 2~m wide. A and C show a 1~m door.\\
\textbf{6.1} 1~mm stands for 5~cm.\\
\textbf{6.2} 1~cm.\\
\textbf{6.3} No. The chair is bigger than the gap.\\
\textbf{6.4} Remeasure the real gap and the model, then see whether a bigger scale is needed.\\
\textbf{6.5} The old model still answers big questions, and it shows what was tried.\\
\textbf{6.6} (b).\\
\textbf{6.7} (a) M (b) R (c) M.\\
\textbf{6.8} No. The model has no door swing. It would need to keep the swing.\\
\textbf{6.9} The chair and gap are at different scales, so the sizes cannot be compared.\\
\textbf{6.10} $5.9\times10=59$, $5.8\times10=58$, difference 1~cm.\\
\textbf{6.11} ``The $1:50$ model could not read a 1~cm margin, and a $1:10$ corner model fixed it.''\\
\textbf{6.12} Test against the real thing, use one scale, or name what was dropped.\\
\textbf{6.13} Without a test, nobody knows which questions the model can answer.\\
\textbf{6.14 to 6.16} Answers vary. Check that each circle's number is accounted for.\\
\textbf{Slop check.} (a) You cannot tell: no title, no labels. (b) Many colours, stars and joke text. (c) A title, a scale, labels for doors. (d) The stars and jokes.\\
\textbf{Whole course.} Answers vary. Sample: Book~2 kept one trial's result and set aside the other trials; Book~4 kept direction and size and set aside the cause of the push.

\clearpage\label{notes}
{\sans\bfseries\LARGE\color{ink} Words}\par
\small
\begin{multicols}{2}
\word{diagram}{a drawing that shows how parts connect, not how they look.}
\word{gap}{the difference between what a model predicts and what is real.}
\word{map}{a flat drawing of a place that keeps places and distances.}
\word{margin}{the spare room a question has, such as a gap minus a chair's width.}
\word{model}{a simple stand-in that keeps some things and drops others.}
\word{prediction}{an answer written down before checking.}
\word{scale}{how big the model is compared with the real thing, such as $1:50$.}
\word{scale bar}{a short ruler on a map that shows the scale.}
\end{multicols}

\vspace{0.4em}
{\sans\bfseries\large\color{sky!70!black} Notes for grown-ups}\par
This book teaches that a model is trustworthy for a question when its predictions about that question were compared with the real thing, and when it states which questions it cannot answer. Each chapter has one activity, one worked example with made-up practice numbers, one trust box, one ``what did we keep?'' prompt and a page of puzzles.

\textbf{Materials and safety.} A tape measure, a ruler marked in millimetres, paper, cardboard, scissors and tape. Use blunt scissors. Measure only rooms and paths that are safe to walk. Stay away from roads, stairs and ladders.

\textbf{Scales.} A scale such as $1:50$ means every length on the model is the real length divided by 50. Children practise dividing and multiplying by 50, 20 and 10, and see that the choice of scale decides what can be shown.

\textbf{The practice numbers} were made up so the arithmetic is easy to check. Real rooms and hallways will give different numbers.

\subhead{Sources}
\begin{itemize}[leftmargin=1.3em,itemsep=0.2em]
\item E. R. Tufte. \emph{Envisioning Information}. Graphics Press, 1990. For what a drawing keeps and what it drops.
\item M. Monmonier. \emph{How to Lie with Maps}, 3rd edition. University of Chicago Press, 2018. For scale and what maps leave out.
\item J. R. Taylor. \emph{An Introduction to Error Analysis}, 2nd edition. University Science Books, 1997. For measurement error and gaps.
\end{itemize}

\clearpage
\thispagestyle{blank}
\begin{tikzpicture}[remember picture,overlay]
  \fill[sun!35] ($(current page.south west)+(0.4in,0.4in)$) rectangle ($(current page.north east)-(0.4in,0.4in)$);
  \node[anchor=north,align=center,font=\sffamily\bfseries,text=ink,text width=5.4in] at ($(current page.north)+(0,-1.4in)$) {\Large That is the Science series\\[6pt]\fontsize{26}{30}\selectfont Nine books, one habit};
  \node[align=center,font=\sffamily,text width=4.6in,text=ink] at (current page.center) {Measure, test fairly, keep a ledger, sort with care and build models you check. Say what you kept and what you set aside. Take these habits into Math, Ecology, Spelling and Grammar, and Problem Solving.};
  \node at ($(current page.south)+(0,2.6in)$) {\scalebox{3}{\multixbot[ok]}};
  \node[font=\sffamily\bfseries,text=ink] at ($(current page.south)+(0,0.9in)$) {Ages 9 to 11 \quad$\star$\quad Science Book 9};
\end{tikzpicture}
